\appendix

\section{How far a repair must travel: an exact bounded-distance screen}\label{app:distance}

The method studied in this paper is a repair: from a colouring that is not yet valid, change edges
until no forbidden clique remains. Whether that is possible at all is a question about distance, and
inside a Hamming ball the question is decidable exactly. This appendix records the answers for six
smaller cells, obtained during the screening stage described in \S\ref{sec:ai}. They are the reason
the present cell was the one worth pushing, and they delimit the method in a way that no search can.

\subsection*{The question}

Let $B$ be a colouring of $K_{n_0+1}$ obtained by adding one vertex to a valid $(r,s)$-colouring of
$K_{n_0}$ --- a twin, in the sense of \S\ref{sec:method}. For $k=0,1,2,\dots$ ask:

\begin{center}
\emph{is there an $(r,s)$-colouring of $K_{n_0+1}$ at Hamming distance at most $k$ from $B$?}
\end{center}

An unsatisfiable instance at bound $k$ is a proof that \emph{no} sequence of at most $k$ edge changes
repairs $B$, whatever heuristic is applied. It is therefore a statement about the cell, not about a
program.

\subsection*{The encoding}

One Boolean variable per edge, with $1$ meaning red. The constraints are that no red $K_r$ and no blue
$K_s$ exist. The red side is usually small enough to write out in full; the blue side is not --- for
$n=40$ there are $\binom{40}{10}=8.5\times10^{8}$ ten-sets --- so the clique constraints are added
lazily: solve, enumerate every violated clique exactly, add the clause forbidding that clique, and
re-solve. The distance bound is imposed by a totalizer over flip variables
$d_e \leftrightarrow (x_e \neq B_e)$, asserting $\sum_e d_e \le k$.

Two solvers were used, CaDiCaL~1.5.3 and MiniSat~2.2, and they agree on every bound that both
completed.

\subsection*{A positive control is mandatory}

An encoding that is too strong returns ``unsatisfiable'' for every input, which looks identical to a
correct proof. The screen was therefore validated against a case whose answer is known: take a valid
$(4,6)$-colouring of $K_{35}$, flip $5$ of its edges at random, and require the encoder to return
satisfiable at $k=5$ with a witness that re-verifies to zero conflicts, and unsatisfiable at $k=4$.
Both behaviours were observed, and the returned witness was checked by an independent enumerator.
Without such a control a distance bound is worthless.

\subsection*{Results}

The column $k^\ast$ is the largest $k$ for which unsatisfiability was proved, so a repair from that
state needs at least $k^\ast+1$ edge changes.

\begin{center}
\small
\begin{tabular}{@{}llrrl@{}}
\toprule
cell & best twin & conflicts after $+1$ & $k^\ast$ & a repair needs\\
\midrule
$(5,6)$ & $v_{56}$, red & $136$ & $16$ & $\ge 17$ edge changes\\
$(4,6)$ & $v_{31}$, red & $24$ & $16$ & $\ge 17$\\
$(3,10)$ & $v_{35}$, red & $6$ & $12$ & $\ge 13$\\
$(4,7)$ & $v_{15}$, red & $40$ & $12$ & $\ge 13$\\
$(3,11)$ & $v_{20}$, red & $8$ & $10$ & $\ge 11$\\
$(4,9)$ & $v_{0}$, red & $72$ & $8$ & $\ge 9$\\
\midrule
$(6,8)$, $(6,9)$, $(8,10)$ & --- & $338$, $2\,830$, $29\,814$ & none & undecided\\
\bottomrule
\end{tabular}
\end{center}

For $(3,10)$ a second, better starting state was also screened: a $40$-vertex colouring with four red
triangles and no blue $K_{10}$, obtained by choosing the new vertex's red neighbourhood rather than
copying a vertex. No valid colouring lies within $10$ edge changes of it.

\subsection*{Two conclusions}

\emph{The method cannot move those six cells from those bases.} For $(4,6)$ this is decisive: every
engine in the study, including one running the complete pipeline of \cite{cremer} with preparation,
risk chains and mutation, stalls at $12$ conflicts, and the nearest valid $36$-vertex colouring is at
least $17$ edge changes away. No amount of tuning, restarting or extra seeds closes a gap of $17$.

\emph{Conflict count does not rank repairability.} The $(4,9)$ start has $72$ conflicts and needs only
$9$ edge changes; the $(3,10)$ start has $6$ --- the smallest in the table, and the natural first
choice --- and needs $13$. Ranking these cells by how many conflicts the vertex addition creates, which
is the obvious heuristic and the one adopted at the outset of this work, orders them wrongly. The
$(8,10)$ cell, which looked the least promising of all with $29\,814$ initial conflicts, is the one
that moved.

\subsection*{Limits}

The bounds are for the specific best-twin extension of each base colouring; they are not statements
about the cells, and a colouring of the target order might exist far from those particular states. For
$(6,8)$, $(6,9)$ and $(8,10)$ no bound was completed --- no rung finishes at those sizes --- so for
those cells nothing is claimed in either direction, including for the result of this paper.
